\chapter{Standard Atmosphere (1976)}

The 1976 US Standard Atmosphere is a reference model that defines how air properties --- temperature, pressure, and density --- change with altitude from sea level up to approximately 85 km. The atmosphere is divided into seven stacked layers, each with a fixed base temperature and a constant lapse rate (or zero for isothermal layers). This model is widely used as a standard reference baseline for aircraft performance analysis, propulsion calculations, and aerodynamic design, providing consistent, internationally agreed-upon values of air properties at any altitude. The calculations and layer constants are implemented in reference to the official U.S. Standard Atmosphere model \cite{std_atmos1976}.

\section{Nomenclature}

\begin{center}
\begin{tabular}{@{}cl@{}}
\toprule
Symbol & Description \\
\midrule
$h$ & Geopotential altitude (m) \\
$P$ & Static pressure (Pa) \\
$\rho$ & Static density (kg/m$^3$) \\
$T$ & Static temperature (K) \\
$a$ & Speed of sound (m/s) \\
$\mu$ & Dynamic viscosity (Pa$\cdot$s) \\
$L$ & Layer lapse rate (K/m), zero for isothermal layers \\
$g_0$ & Sea-level gravitational acceleration (m/s$^2$) \\
$R$ & Specific gas constant of air (J/(kg$\cdot$K)) \\
$\gamma$ & Specific heat ratio of air \\
$\mu_0$ & Sutherland reference dynamic viscosity (Pa$\cdot$s) \\
$T_{Suth}$ & Sutherland reference temperature (K) \\
$S_{Suth}$ & Sutherland temperature constant (K) \\
\bottomrule
\end{tabular}
\end{center}

\subsection*{Constants}

\begin{center}
\begin{tabular}{@{}cl@{}}
\toprule
Symbol & Value \\
\midrule
$g_0$ & $9.80665$ m/s$^2$ \\
$R$ & $287.0528$ J/(kg$\cdot$K) \\
$\gamma$ & $1.4$ \\
$\mu_0$ & $1.716\times10^{-5}$ Pa$\cdot$s (Sutherland reference) \\
$T_{Suth}$ & $273.15$ K \\
$S_{Suth}$ & $110.4$ K \\
\bottomrule
\end{tabular}
\end{center}

\section{Governing Idea}

The 1976 US Standard Atmosphere models the atmosphere as seven stacked
layers from sea level to 84{,}852 m, each with a fixed base temperature
and a constant lapse rate (zero for isothermal layers). Within a layer,
temperature is linear in altitude and pressure follows in closed form
from the hydrostatic equation, so the model is exact and non-iterative:
whichever quantity the user supplies, the containing layer is located
first, and the layer equation is then evaluated (given altitude) or
inverted in closed form (given pressure or density) to recover the full
atmospheric state. Density, speed of sound, and viscosity always follow
from the ideal-gas law, the isentropic sound-speed relation, and
Sutherland's law respectively, once $T$ and $P$ are known.

\section{Layer Table}

\begin{center}
\begin{tabular}{@{}cccc@{}}
\toprule
$h_{base}$ (m) & $T_{base}$ (K) & $L$ (K/m) & Layer \\
\midrule
0 & 288.15 & $-0.0065$ & Troposphere \\
11{,}000 & 216.65 & 0 & Tropopause \\
20{,}000 & 216.65 & $+0.001$ & Stratosphere I \\
32{,}000 & 228.65 & $+0.0028$ & Stratosphere II \\
47{,}000 & 270.65 & 0 & Stratopause \\
51{,}000 & 270.65 & $-0.0028$ & Mesosphere I \\
71{,}000 & 214.65 & $-0.002$ & Mesosphere II \\
84{,}852 & 186.946 & --- & Top boundary \\
\bottomrule
\end{tabular}
\end{center}

For a lapse-rate layer ($L\neq0$), the temperature and pressure at the top
of the layer follow from

\begin{equation}
T_{top} = T_{base} + L(h_{top}-h_{base}), \qquad
P_{top} = P_{base}\left(\frac{T_{top}}{T_{base}}\right)^{-\frac{g_0}{LR}}
\end{equation}

and for an isothermal layer ($L=0$),

\begin{equation}
T_{top} = T_{base}, \qquad
P_{top} = P_{base}\exp\left(-\frac{g_0(h_{top}-h_{base})}{RT_{base}}\right)
\end{equation}

The base pressure of each layer is the top pressure of the one before it.
Layer density, where needed for locating a layer by density, follows from
$\rho_{top} = P_{top}/(RT_{top})$.

\section{Locating the Layer}

\begin{itemize}[nosep]
    \item \textbf{By altitude:} the layer whose $[h_{base}, h_{top}]$
    interval contains $h$.
    \item \textbf{By pressure:} pressure decreases monotonically with
    altitude, so the layers are scanned from the ground up and the first
    layer whose top-of-layer pressure is at or below the given $P$ is
    selected.
    \item \textbf{By density:} the same monotonic scan, but comparing
    against each layer's top-of-layer density.
\end{itemize}

\section{Calculation Cases}

Depending on the input parameter provided by the user, the standard atmosphere properties are computed using one of the three cases described below.

\subsection*{Case 1: Given Altitude ($h$)}

Valid for $0 \le h \le 84{,}852$ m. With the containing layer identified:

\begin{equation}
T = T_{base} + L(h-h_{base}), \qquad
P = P_{base}\left(\frac{T}{T_{base}}\right)^{-\frac{g_0}{LR}} \qquad (L\neq0)
\end{equation}

\begin{equation}
T = T_{base}, \qquad
P = P_{base}\exp\left(-\frac{g_0(h-h_{base})}{RT_{base}}\right) \qquad (L=0)
\end{equation}

\begin{equation}
\rho = \frac{P}{RT}
\end{equation}

\subsection*{Case 2: Given Static Pressure ($P$)}

Valid for $P_{min} < P \le P_{max}$, where $P_{max}$ is the sea-level base
pressure and $P_{min}$ is the pressure at the top of the highest layer.
With the containing layer identified, the Case~1 relation is inverted in
closed form:

\begin{equation}
T = T_{base}\left(\frac{P}{P_{base}}\right)^{-\frac{LR}{g_0}}, \qquad
h = h_{base} + \frac{T-T_{base}}{L} \qquad (L\neq0)
\end{equation}

\begin{equation}
T = T_{base}, \qquad
h = h_{base} - \frac{RT_{base}}{g_0}\ln\left(\frac{P}{P_{base}}\right) \qquad (L=0)
\end{equation}

\begin{equation}
\rho = \frac{P}{RT}
\end{equation}

\subsection*{Case 3: Given Static Density ($\rho$)}

Valid for $\rho_{min} < \rho \le \rho_{max}$, where $\rho_{max}$ is the
sea-level density and $\rho_{min}$ is the density at the top of the
highest layer. With the containing layer identified (base density
$\rho_{base} = P_{base}/(RT_{base})$), the density-altitude relation is
inverted in closed form:

\begin{equation}
T = T_{base}\left(\frac{\rho}{\rho_{base}}\right)^{-\frac{LR}{g_0+LR}}, \qquad
h = h_{base} + \frac{T-T_{base}}{L} \qquad (L\neq0)
\end{equation}

\begin{equation}
T = T_{base}, \qquad
h = h_{base} - \frac{RT_{base}}{g_0}\ln\left(\frac{\rho}{\rho_{base}}\right) \qquad (L=0)
\end{equation}

\begin{equation}
P = \rho RT
\end{equation}

\section{Derived Properties}

Once $T$ is known by any of the three routes above, the remaining derived outputs follow directly:

\subsection*{Speed of Sound ($a$)}
The local speed of sound in air is calculated using the thermodynamic relation:
\begin{equation}
a = \sqrt{\gamma R T}
\end{equation}

\subsection*{Dynamic Viscosity ($\mu$)}
The dynamic viscosity of air is modeled as a function of temperature using Sutherland's law:
\begin{equation}
\mu = \mu_0\left(\frac{T}{T_{Suth}}\right)^{3/2}\frac{T_{Suth}+S_{Suth}}{T+S_{Suth}}
\end{equation}

\section{Program Flow and Architecture}

The Standard Atmosphere module accepts exactly one of three inputs and produces all atmospheric outputs via closed-form ISA layer equations.

\begin{center}
\begin{tikzpicture}[x=1cm,y=-1cm,
    inp/.style={rectangle, draw=blue!70!black, thick, fill=blue!6, align=center,
                text width=3.2cm, minimum height=0.7cm, font=\small\sffamily},
    eng/.style={rectangle, draw=orange!80!black, thick, fill=orange!8, align=center,
                text width=3.8cm, font=\small\sffamily},
    outn/.style={rectangle, draw=green!60!black, thick, fill=green!6, align=center,
                text width=2.8cm, minimum height=0.7cm, font=\small\sffamily},
    arr/.style={-{Stealth[scale=1.0]}, thick, draw=black!70},
    bus/.style={thick, draw=orange!70!black}
]
%--- Inputs (y increases downward) ---
\node[inp] (h)   at (0, 0.4) {Case 1: Altitude $h$};
\node[inp] (Pn)  at (0, 1.4) {Case 2: Pressure $P$};
\node[inp] (rho) at (0, 2.4) {Case 3: Density $\rho$};
%--- Engine ---
\node[eng, minimum height=2.4cm] (eng) at (4.8, 1.4)
    {\textbf{AtmosphereEngine}\\[3pt]
    Layer lookup\\[-1pt]
    Closed-form solve\\[-1pt]
    Sutherland viscosity};
%--- Outputs ---
\node[outn] (T1) at (9.4, 0.0) {Temperature $T$};
\node[outn] (P1) at (9.4, 0.9) {Pressure $P$};
\node[outn] (R1) at (9.4, 1.8) {Density $\rho$};
\node[outn] (A1) at (9.4, 2.7) {Speed of sound $a$};
\node[outn] (M1) at (9.4, 3.6) {Viscosity $\mu$};
%--- Input arrows ---
\draw[arr] (h.east)   -- (eng.west |- h.east);
\draw[arr] (Pn.east)  -- (eng.west);
\draw[arr] (rho.east) -- (eng.west |- rho.east);
%--- Bus on right of engine ---
\coordinate (busX) at ($(eng.east)+(0.55,0)$);
\draw[bus] (eng.east) -- (busX);
\draw[bus] (busX |- T1) -- (busX |- M1);
%--- Output arrows ---
\draw[arr] (busX |- T1) -- (T1.west);
\draw[arr] (busX |- P1) -- (P1.west);
\draw[arr] (busX |- R1) -- (R1.west);
\draw[arr] (busX |- A1) -- (A1.west);
\draw[arr] (busX |- M1) -- (M1.west);
\end{tikzpicture}
\end{center}

\noindent The user provides one of $h$, $P$, or $\rho$. The engine finds the correct ISA layer, inverts the layer equations analytically for temperature and altitude, then computes $a=\sqrt{\gamma R T}$ and $\mu$ via Sutherland's law. All paths are closed-form; no iteration is required.
